\newpage
\section{Projekt}

Niniejszy rozdział opisuje decyzje projektowe podjęte w trakcie realizacji systemu,
obejmując wybór technologii komunikacji radiowej i platformy sprzętowej, konfigurację
sieci mesh oraz architekturę aplikacji mobilnej.

\subsection*{Wybór technologii}

Dobór platformy sprzętowej poprzedzono wyborem techniki modulacji. W projekcie wykorzystano technologię
LoRa. Poniżej przedstawiono krótki opis teoretyczny oraz uzasadnienie wyboru.

\subsection*{Modulacja LoRa i technika Chirp Spread Spectrum}

LoRa (ang.\ \textit{Long Range}) to technika modulacji radiowej, która pozwala na przesyłanie
niewielkich pakietów danych na bardzo duże odległości - nawet kilkanaście kilometrów - przy
minimalnym zużyciu energii.\footnote{\url{https://www.gov.pl/web/instytut-lacznosci/lora-pod-lupa}}
Jej zamysł opiera się na technice modulacji CSS (ang.\ \textit{Chirp Spread Spectrum}) lub FSCM
(ang.\ \textit{Frequency Shift Chirp Modulation})~\cite{vangelista2017lora}. Polega ona na
kodowaniu sygnału przez cykliczne przesunięcie punktu startowego chirpa. Wspomniany chirp to
,,blok'' sygnału o liniowo zmiennej częstotliwości.

Modulację LoRa opisują trzy główne parametry:
\begin{itemize}
    \item \textbf{Współczynnik rozpraszania SF} (ang.\ \textit{Spreading Factor}) - określa
          liczbę chirpów przypadających na jeden symbol. Im wyższy SF, tym dłuższy czas transmisji
          symbolu, ale też lepsza czułość odbiornika i większy zasięg. W praktyce stosuje się
          wartości od 7 do 12.

    \item \textbf{Szerokość pasma BW} (ang.\ \textit{Bandwidth}) - określa zakres częstotliwości
          zajmowany przez sygnał. Szersza szerokość pasma oznacza krótszy czas transmisji symbolu
          przy tym samym SF, ale też wyższy poziom szumów i mniejszy zasięg. W praktyce do wyboru
          są częstotliwości 125, 250 lub 500~kHz.

    \item \textbf{Współczynnik kodowania CR} (ang.\ \textit{Coding Rate}) - określa nadmiarowość
          kodu korekcji błędów. Wyższy CR zwiększa odporność na błędy transmisji kosztem wydajności.
\end{itemize}

Przepustowość bitową modulacji LoRa wyraża wzór:\footnote{\url{https://josefmtd.com/2018/08/14/spreading-factor-bandwidth-coding-rate-and-bit-rate-in-lora-english/}}

\begin{equation}
    R_b = SF \cdot \frac{BW}{2^{SF}} \cdot \frac{CR}{4+CR}
    \label{eq:lora-bitrate}
\end{equation}

Kluczową zaletą CSS jest odporność na zakłócenia wąskopasmowe poprzez rozłożenie energii sygnału
na całe dostępne pasmo. Efekt Dopplera (istotny przy komunikacji w ruchu) również jest ograniczony.
Wynika to z tego, że odczyt symbolu polega w skrócie na mnożeniu odebranego sygnału przez chirp
bazowy i wykonaniu FFT (ang.\ \textit{Fast Fourier Transform}). Powstały w wyniku pik widmowy jest
jedynie przesunięty i przy prędkościach spotykanych w praktyce nie spowoduje błędu dekodowania.

\subsection*{Wybór platformy sprzętowej}

W projekcie rozważono możliwość wykorzystania różnych platform sprzętowych z modułami radiowymi LoRa:
\begin{enumerate}
    \item \textbf{Raspberry Pi z modułami LoRa} - mimo licznych zalet, takich jak bogate wsparcie,
          względnie wysoka moc obliczeniowa oraz elastyczność, przeważyły wady, które wykluczyły to
          rozwiązanie: zbyt duże zużycie energii, duże rozmiary, brak wodoszczelności. Podobne
          argumenty zdecydowały o wyeliminowaniu platformy \textit{Arduino~Uno/Nano z modułem LoRa}.
    \item \textbf{ESP32 z modułami LoRa} - w szczególności rozważono \textit{Heltec ESP32 LoRa},
          czyli płytkę z wbudowanym modułem LoRa i wyświetlaczem OLED. Produkt ten posiada wiele
          cech wypełniających wymagania projektu (wbudowany moduł Bluetooth i LoRa, niskie zużycie
          energii, niski koszt), jednak o odrzuceniu tej platformy zadecydowały pragmatyczne względy: 
          brak wodoszczelności - ewentualna konieczność druku 3D obudowy, a osiądnięcie IP65 w drukowanej 
          obudowie wymaga m.in. dedykowanych złączy kablowych, co znacznąco podnosi poziom skomplikowania projektu.
    \item \textbf{Meshtastic na dedykowanym sprzęcie LoRa} - otwartoźródłowa platforma łącząca
          firmware dla modułów radiowych LoRa z ekosystemem aplikacji klienckich. Tworzy
          zdecentralizowaną sieć mesh w paśmie ISM bez infrastruktury zewnętrznej, każdy węzeł
          automatycznie retransmituje pakiety pozostałych, rozszerzając zasięg sieci wraz
          z liczbą urządzeń. Platforma udostępnia publiczne API w postaci aplikacji Android
          \texttt{com.geeksville.mesh}\footnote{Repozytorium aplikacji Meshtastic na Android:
          \url{https://github.com/meshtastic/Meshtastic-Android}. Dokumentacja API:
          \url{https://meshtastic.org/docs/development/firmware/build/}.} oraz eksportuje 
          interfejs AIDL\footnote{AIDL (ang.\ \textit{Android Interface Definition Language}) to mechanizm
          systemu Android umożliwiający komunikację między procesami.} umożliwiający aplikacjom
          trzecim odczyt danych telemetrycznych - w tym pozycji GPS węzłów.
          Sprzęt kompatybilny z Meshtastic jest ogólnodostępny i tani, a samo oprogramowanie
          bezpłatne. Spośród rozważanych opcji platforma ta najlepiej spełnia wymagania projektu.
\end{enumerate}

Wybór platformy Meshtastic podyktowany był przede wszystkim dostępnością dedykowanego sprzętu
spełniającego wymagania wodoszczelności oraz możliwością integracji z aplikacją mobilną przez
publiczne API. W ramach platformy wybór padł na:
\begin{itemize}
    \item \textit{T-Beam} - popularna platforma oparta na ESP32 z wbudowanym modułem GPS, anteną
          LoRa oraz akumulatorem. Posiada duży akumulator oraz otwarte piny, które zostawiają
          szerokie pole manewru do rozwoju tego węzła w roli stacji bazowej.
    \item \textit{T-2000-E} - kompaktowe urządzenie oferujące wodoszczelność (IP65), wbudowany GPS
          oraz zoptymalizowaną żywotność baterii. Wymiary i kształt są spełniają wymagania projektu, bo
          urządzenie ma wymiary dwóch złożonych kart bankowych - 6,5x85x55mm i waży 32 g.
\end{itemize}

Obie platformy przedstawia rysunek~\ref{fig:platformy}.

\begin{figure}[H]
    \centering
    \begin{subfigure}{0.48\linewidth}
        \centering
        \includegraphics[width=\linewidth]{tbeam.png}
        \caption{T-Beam - węzeł stacji bazowej}
        \label{fig:tbeam}
    \end{subfigure}
    \hfill
    \begin{subfigure}{0.48\linewidth}
        \centering
        \includegraphics[width=\linewidth]{tracker-t1000-e.png}
        \caption{T1000-E - węzeł obroży}
        \label{fig:t1000e}
    \end{subfigure}
    \caption{Wybrane platformy sprzętowe Meshtastic}
    \label{fig:platformy}
\end{figure}

Otwartoźródłowy charakter projektu stwarza możliwość implementacji rozszerzeń dla węzłów sieci
w postaci modułów. Dzięki temu wybór Meshtastic jako platformy sprzętowej nie ograniczy projektu
do domyślnie dostępnych funkcji. Spośród dostępnych konfiguracji radiowych można wybrać parametry
dla modulacji LoRa. Zostały one zestawione w tabeli~\ref{tab:lora-presets}.\footnote{\url{https://meshtastic.org/docs/software/site-planner/}}

\begin{table}[!h]
    \caption{Predefiniowane ustawienia kanałów LoRa w Meshtastic}
    \label{tab:lora-presets}
    \centering
    \begin{tabular}{|l|c|c|c|}
        \hline
        \textbf{Ustawienie} & \textbf{BW [kHz]} & \textbf{SF} & \textbf{CR} \\
        \hline
        ShortTurbo   & 500 & 7  & 4/5 \\
        ShortFast    & 250 & 7  & 4/5 \\
        ShortSlow    & 250 & 8  & 4/5 \\
        MediumFast   & 250 & 9  & 4/5 \\
        MediumSlow   & 250 & 10 & 4/5 \\
        LongFast     & 125 & 11 & 4/5 \\
        LongModerate & 125 & 11 & 4/8 \\
        LongSlow     & 125 & 12 & 4/8 \\
        \hline
    \end{tabular}
\end{table}


Zazwyczaj do konfiguracji sieci Meshtastic używa się następujących parametrów: $SF=12$,
$BW=125$~kHz oraz $CR=4/8$. Podstawiając je do wzoru~(\ref{eq:lora-bitrate}) otrzymuje się przepustowość
na poziomie 250~bps. Jest ona w zupełności wystarczająca do przesyłania krótkich pakietów
zawierających współrzędne GPS.

\subsection*{Konfiguracja sieci Meshtastic}

Sieć Meshtastic skonfigurowano tak, aby wszystkie węzły mogły przesyłać między sobą informacje
o swoim położeniu GPS. W sieci Meshtastic informacje o pozycji GPS są przesyłane jako część telemetrii węzła. Telemetria
jest automatycznie udostępniana wszystkim węzłom w sieci, które używają tego samego kanału PRIMARY
(lub kanału wtórnego z włączoną pozycją), mają zgodne ustawienia LoRa (częstotliwość, modem preset)
oraz są w zasięgu (bezpośrednio lub przez pośrednie węzły). Schemat sieci przedstawia rysunek~\ref{fig:schematSieci}.

\begin{figure}[H]
    \label{fig:schematSieci}
    \centering
    \includegraphics[width=0.9\linewidth]{schematSieci.png}
    \caption{Schemat sieci Meshtastic z pięcioma węzłami}
\end{figure}

\paragraph{Konfiguracja - prywatne kanały.}

Aby udostępniać pozycję wyłącznie węzłom należącym do sieci prywatnej, konieczne było zastosowanie
prywatnego kanału. Konfiguracja przebiegła następująco:

\begin{enumerate}
    \item \textbf{Wyłączono pozycję na kanale PRIMARY}:
          Settings $\rightarrow$ Radio Config $\rightarrow$ Channels $\rightarrow$ Channel 0 (PRIMARY)
          $\rightarrow$ Position Sharing: \textbf{WYŁĄCZONE}.

    \item \textbf{Utworzono prywatny kanał}:
          Settings $\rightarrow$ Radio Config $\rightarrow$ Channels $\rightarrow$ \textbf{Add Channel},
          z następującymi parametrami:
          \begin{itemize}
              \item \textbf{Channel Name}: \texttt{Reja}
              \item \textbf{PSK}: wspólny klucz szyfrowania
              \item \textbf{Position Sharing}: \textbf{WŁĄCZONE}
              \item \textbf{Position Precision}: najwyższy poziom precyzji 32-bity
                    (deklaruje dokładność co do metra)
          \end{itemize}

    \item \textbf{Dodano ten sam kanał do wszystkich węzłów w grupie}: wszystkie węzły mają
          identyczną nazwę kanału \texttt{Reja} oraz identyczny klucz PSK.
\end{enumerate}


